\documentclass[11pt]{amsart}

\usepackage{amsmath}
\usepackage{amssymb}
\usepackage{amsthm}
%\usepackage{psfig}

\begin{document}
\baselineskip=15pt
\textheight=8.4in
\parindent=0pt
\def\sk {\hskip .5cm}
\def\skv {\vskip .12cm}
\def\cos {\mbox{cos}}
\def\sin {\mbox{sin}}
\def\tan {\mbox{tan}}
\def\intl{\int\limits}
\def\lm{\lim\limits}
\newcommand{\frc}{\displaystyle\frac}
\def\xbf{{\mathbf x}}
\def\fbf{{\mathbf f}}
\def\gbf{{\mathbf g}}

\def\Ker{{\rm Ker\,}}
\def\phi{\varphi}

\def\dbA{{\mathbb A}}
\def\dbB{{\mathbb B}}
\def\dbC{{\mathbb C}}
\def\dbD{{\mathbb D}}
\def\dbE{{\mathbb E}}
\def\dbF{{\mathbb F}}
\def\dbG{{\mathbb G}}
\def\dbH{{\mathbb H}}
\def\dbI{{\mathbb I}}
\def\dbJ{{\mathbb J}}
\def\dbK{{\mathbb K}}
\def\dbL{{\mathbb L}}
\def\dbM{{\mathbb M}}
\def\dbN{{\mathbb N}}
\def\dbO{{\mathbb O}}
\def\dbP{{\mathbb P}}
\def\dbQ{{\mathbb Q}}
\def\dbR{{\mathbb R}}
\def\dbS{{\mathbb S}}
\def\dbT{{\mathbb T}}
\def\dbU{{\mathbb U}}
\def\dbV{{\mathbb V}}
\def\dbW{{\mathbb W}}
\def\dbX{{\mathbb X}}
\def\dbY{{\mathbb Y}}
\def\dbZ{{\mathbb Z}}


\def\char{{\rm char}}
\def\wt{{\rm wt}}
\def\Aut{{\rm Aut}}
\def\Hom{{\rm Hom}}
\def\End{{\rm End}}
\def\Irr{\mathrm {Irr}}
\def\Im{{\rm Im}}
\def\GL{{\rm GL}}
\def\SL{{\rm SL}}
\def\SNF{{\rm SNF}}
\def\tr {{\rm tr}}
\def\rk{{\rm rk}}
\def\lam{{\lambda}}

\def\la{{\langle}}
\def\ra{{\rangle}}

\bf\centerline{Homework \# 8, to be submitted on Canvas by 11:59pm on Fri, Mar 27th.}\rm
\vskip .2cm
{\bf Plan for the next 2 classes:} Field extensions, algebraic closures and splitting fields. References: 13.1, 13.2 and 13.4 in Dummit and Foote, Lectures 14-16 from Spring 2010 and Lectures 9-11 from Spring 2021.
\vskip .1cm

{\bf Note on hints:} Some hints are given at the end of the assignment, each on a separate page.
Problems (or parts of problems) for which a hint is available at the end are marked with *.

\skv
{\bf 1.} Given a ring $R$ with $1$, the opposite ring $R^{\rm{op}}$ is defined as follows. As a group with addition, $R^{\rm{op}}$ is the same as $R$, and the multiplication $*$ in $R^{\rm{op}}$ is define by $a*b=ba$ (where $ba$ is the product of $b$ and $a$ in $R$).

\begin{itemize}
\item[(a)] Let $\End_R(R)$ be the ring of endomorphisms of $R$ considered as a left module over itself. Prove that $\End_R(R)$ is isomorphic to $R^{\rm{op}}$.
\item[(b)] Let $G$ a group, $K$ a commutative ring with $1$  and $KG$ the corresponding group ring. Prove that $KG$ is isomorphic to its opposite ring.
\end{itemize}

\skv
{\bf 2.} Let $p$ be a prime and $G={\rm Heis}(\dbF_p)$, the Heisenberg group over $\dbF_p$ defined in HW\#7.7.
\begin{itemize}
\item[(a)] Determine the number of conjugacy classes of $G$ and their sizes. As in HW\#7.7, you can work directly with matrices or 
with their expressions in terms of the generators $x,y,z$ introduced in HW\#7.7.
\item[(b)] Let $\omega\neq 1$ be a $p^{\rm th}$ root of unity, that is, $\omega=e^{\frac{2\pi k i}{p}}$ with $1\leq k\leq p-1$.
Let $V$ be a $p$-dimensional complex vector space with basis $e_{[0]}, e_{[1]},\ldots, e_{[p-1]}$ where we think of indices as elements of $\dbF_p$.
Prove that there exists a representation $(\rho_{\omega},V)$ of $G$ such that
\begin{itemize}
\item $\rho_{\omega}(z) e_{[k]}=\omega e_{[k]}$ for each $k$ (that is, $\rho_{\omega}(z)$ is just the scalar multiplication by $\omega$),
\item $\rho_{\omega}(y) e_{[k]}=e_{[k+1]}$ for each $k$ (that is, $\rho_{\omega}(y)$ cyclically permutes the basis vectors) and finally 
\item $\rho_{\omega}(x) e_{[k]}=\omega^{k} e_{[k]}$ for each $k$.
\end{itemize}
\item[(c)] Prove that every representation in (b) is irreducible (do not do this directly from definition) and every irreducible complex representation of $G$ is either one-dimensional 
or equivalent to $(\rho_{\omega},V)$ for some $\omega$.
\end{itemize}
\skv
{\bf 3.*} Let $G=S_n$ for some $n\geq 2$ and $\chi$ an irreducible complex character of $G$. Prove that $\chi$ is real-valued, that is,
$\chi(g)\in \dbR$ for all $g\in G$. 
\skv
{\bf 4.} Give an example of two representations $V$ and $W$ of the same group which are not equivalent, but have the same dimension and the same character. Recall that by Proposition~17.6 from class this cannot happen if $G$ is finite and representations are complex.
\skv
{\bf 5.} Let $C$ be the cube in $\dbR^3$ whose vertices have coordinates 
$(\pm 1, \pm 1,\pm 1)$. Let $G$ be the group of rotations of $C$, that is rotations in $\dbR^3$ which preserve the cube (you may assume that $G$ is a group without proof). Let $X$ be the set of $4$ main diagonals of $C$ (diagonals
connecting the opposite vertices). Note that $G$ naturally acts on $X$ and therefore we have a homomorphism $\pi:G\to Sym(X)\cong S_4$.
\begin{itemize}
\item[(a)] (practice, no need to turn in) Prove that $\pi$ is an isomorphism.
\item[(b)] Note that $G$ is naturally a subgroup of $\GL_3(\dbR)$ and hence
also a subgroup of $\GL_3(\dbC)$, and let $\iota: G\to \GL_3(\dbC)$
be the inclusion map. By (a) we get a representation $\iota\circ \pi^{-1}:
S_4\to\GL_3(\dbC)$. Prove that this representation is equivalent to the tensor product of the standard and sign representations.
\end{itemize}

\skv
{\bf 6.} The goal of this problem is to explicitly decompose the regular representation $(\rho_{reg},\dbC[S_3])$ as a direct sum of irreducible representations of $S_3$. Recall that
$S_3$ has 3 I$\dbC$R's: two one-dimensional (the trivial representation and the sign representation) and one two-dimensional (the standard representation) and that each I$\dbC$R appears in $\dbC[S_3]$ with multiplicity equal to its dimension (Theorem~17.4 from class).
\begin{itemize}
\item[(a)] Let $H=\la (1,2,3)\ra$ and consider $(\rho_{reg},\dbC[S_3])$ as a representation of $H$. Prove that $\dbC[S_3]=V_1\oplus V_2$ where both 
$V_1$ and $V_2$ are $H$-invariant and equivalent (as $H$-representations) 
to the regular representation of $H$.
\item[(b)]  Now explicitly decompose $\dbC[H]$, the regular representation of $H$, into a direct sum of 3 one-dimensional $H$-representations
({\bf Hint}: Since $H$ is cyclic of order $3$, this boils down to finding eigenspaces for a suitable endomorphism of $\dbC[H]$ of order $3$).
Combining this with (a), we get an explicit decomposition
$\dbC[S_3]=\oplus_{i=1}^6 W_i$ where each $W_i$ is a one-dimensional 
and $H$-invariant. 

Show that after a suitable renumbering of $W_1,\ldots, W_6$
the following is true: $W_1\oplus W_2$ and $W_3\oplus W_4$ are both irreducible
subrepresentations of $S_3$ (these are the copies of the standard representation we are supposed to get by Theorem~17.4), while $W_5\oplus W_6$ decomposes (in a different way) into a direct sum of the trivial and the sign representations of $S_3$. 
\end{itemize}
\skv

\newpage
{\bf Hint for 3:} Use our discussion in Lecture~15 to find a general condition on a finite group $G$ which guarantees that all of its complex characters are real-valued and then show that this condition holds for $S_n$.
\end{document}
